\documentclass[11pt,letterpaper]{article}
\usepackage[T1]{fontenc}
\usepackage[margin=0.85in]{geometry}
\usepackage{amsmath,amssymb,booktabs,tabularx,microtype}
\usepackage{xcolor,listings}
\usepackage[hyphens]{url}
\usepackage[colorlinks=true,linkcolor=blue!45!black,urlcolor=blue!45!black,
  citecolor=blue!45!black]{hyperref}
\hypersetup{pdftitle={The VOID-collapse model in 3D Truchas flow},
  pdfsubject={Source-based description of the projection and volume-tracking algorithm}}
\lstset{basicstyle=\small\ttfamily,columns=fullflexible,keepspaces=true,
  breaklines=true,frame=single,rulecolor=\color{black!25},
  xleftmargin=0.4em,xrightmargin=0.4em}
\setlength{\parindent}{0pt}
\setlength{\parskip}{0.5em}
\newcommand{\code}[1]{\texttt{\detokenize{#1}}}
\newcommand{\rev}{344a7c3a98d5046533585e6fedc87ffee312388b}
\newcommand{\srclink}[3]{\href{https://gitlab.com/truchas/truchas/-/blob/344a7c3a98d5046533585e6fedc87ffee312388b/#1\#L#2}{#3}}

\title{The VOID-collapse model in 3D Truchas flow}
\author{Implementation and mathematical interpretation}
\date{September 12, 2026}

\begin{document}
\maketitle

\begin{abstract}
Truchas provides an experimental model for collapsing small regions of VOID
entrained in liquid. It introduces a prescribed negative divergence into the
pressure projection in selected mixed liquid/void cells. The corrected
velocity subsequently drives volume tracking; when a cell becomes overfilled,
the geometric tracker preferentially removes VOID. This note derives the
implemented divergence source, explains cell selection and the time-level
coupling, and distinguishes the algorithm from a prescribed decay law or a
physical gas-compression model.
\end{abstract}

\textbf{Source basis.} This describes the refactored, mainline 3D flow code at
commit \code{344a7c3a} (full revision \texttt{\rev}). The relevant sources are
also identical in the 2D worktree at \code{ba178536}. It does not describe
the older v3.0.1 flow implementation or an implementation of collapse in 2D.
Equations below are algebraic consequences of the source unless explicitly
identified as illustrative assumptions. No simulation was run for this note.

\section{Purpose and physical interpretation}

The reference manual explains the motivation: volume tracking can leave
small fragments of VOID behaving effectively as incompressible bubbles,
which persist in the liquid. Enabling \code{void_collapse} allows a
compressive velocity field to develop around these fragments. The manual
labels both the model and its relaxation parameter experimental~\cite{manual}.

VOID represents empty space. This model contains no gas equation of state,
gas mass, gas energy, bubble pressure law, or surface-tension force. In
particular, it does not solve a physical bubble-radius evolution equation.
The new term is an empirical modification of the divergence constraint,
coupled to the tracker's special treatment of VOID.

Pure void cells already provide a zero-pressure condition to the ordinary
free-surface projection. Collapse adds a mechanism for mixed cells without
an immediately adjacent pure void cell. Enabling collapse does not directly
delete volume fractions in the projection routine~\cite{projection,tracker}.

\section{Volume fractions and cell selection}

Let cell $c$ have fixed geometric volume $V_c$. Define
\begin{align}
 F_c &= \sum_{\text{real fluid phases and VOID}} f_{m,c}, &
 L_c &= \sum_{\text{real fluid phases}} f_{m,c},\\
 \alpha_c &= F_c-L_c, & S_c &= 1-F_c.
\end{align}
Here $\alpha_c$ is the VOID fraction of the \emph{whole cell}; it is not
normalized by the fluid-accessible fraction $F_c$. Non-fluid phases are
represented by the lumped solid fraction $S_c$.

\begin{center}
\begin{tabularx}{\linewidth}{@{}l l X@{}}
\toprule
Symbol & Source component & Meaning\\
\midrule
$F_c$ & \texttt{props\%vof(c)} & Fluid fraction including VOID\\
$L_c$ & \texttt{props\%vof\_novoid(c)} & Real-fluid fraction excluding VOID\\
$\alpha_c^n$ & \code{vof_n - vof_novoid_n} & VOID fraction at the last accepted state\\
$\delta\alpha_c$ & \texttt{props\%void\_delta\_cc(c)} & Current fraction minus last accepted fraction\\
$\lambda$ & \code{void_collapse_relaxation} & Dimensionless relaxation factor\\
\bottomrule
\end{tabularx}
\end{center}

The executable classification in \texttt{flow\_props\%update\_cc} is as
follows~\cite{props}. The order of these tests matters.
\begin{enumerate}
\item If $F_c<\varepsilon_f$, classify the cell as \code{solid_t}.
\item Otherwise, if $L_c=0$, classify it as \code{void_t}.
\item Otherwise, if $F_c>L_c$, provisionally classify it as
  \code{regular_void_t}.
\item Otherwise classify it as \code{regular_t}.
\end{enumerate}
The internal \code{cutvof} default is $\varepsilon_f=0.01$. These are the
actual tests in \code{update_cc}; the introductory description in
\code{flow_domain_types.F90} does not fully describe them.

A second pass changes a provisional \code{regular_void_t} cell to
\code{regular_t} if \emph{any face neighbor} is \code{void_t}.
Boundary entries with no neighboring cell are skipped. Thus a cell is
eligible for collapse precisely when it retains the mixed-cell label and
the collapse option is enabled. Both regular cell labels are active in the
pressure solve.

This is a local adjacency test, not a connected-component or bubble-size
test. A collection of mixed cells can be eligible even when its VOID forms
a connected region. Conversely, mixed cells next to a pure void cell are
excluded. Solid neighbors do not themselves exclude a mixed cell, and the
presence of a physical boundary is not a separate exclusion criterion.

\textbf{Cell and face labels have different meanings.} For a cell,
\code{regular_void_t} identifies the selected mixed state just described.
For a face, the same integer label identifies an active-fluid/\code{void_t}
connection, where the pressure on the void side is zero. A collapse-eligible
cell need not have any such face; indeed the neighbor test excludes that
configuration~\cite{props,projection}.

\section{The implemented projection source}

\subsection{History-dependent collapse amount}

The property update computes
\begin{equation}
 \delta\alpha_c = \alpha_c-\alpha_c^n.
 \label{eq:delta}
\end{equation}
This is an Eulerian cellwise change. It can reflect transport and whatever
other updates have changed the current material/phase distribution; it is
not a material derivative or a quantity extracted from the predicted
velocity divergence~\cite{props,cycle}.

For an eligible cell define the nonnegative fraction
\begin{equation}
 q_c=\min\bigl(|\delta\alpha_c|,\lambda\alpha_c\bigr).
 \label{eq:q}
\end{equation}
For other cells, or when collapse is disabled, set $q_c=0$. The mainline
implementation expresses this with two sign branches. After the
predicted-velocity flux sum has been accumulated, it performs the following
operations; the last division acts on the \emph{entire} right-hand side:
\begin{lstlisting}
void = lambda * (props%vof(c) - props%vof_novoid(c))
if (props%void_delta_cc(c) < 0) then
  rhs(c) = rhs(c) + V(c)*max(props%void_delta_cc(c), -void)/dt
end if
if (props%void_delta_cc(c) > 0) then
  rhs(c) = rhs(c) - V(c)*min(props%void_delta_cc(c),  void)/dt
end if
rhs(c) = rhs(c)/dt
\end{lstlisting}
This excerpt abbreviates component names but preserves the signs and
scaling in \code{setup_solver}~\cite{projection}. Both nonzero branches
\emph{subtract} $V_cq_c/\Delta t$ before the final division. Hence the
additional pressure-system right-hand side is
\begin{equation}
 \boxed{\Delta b_c=-\frac{V_cq_c}{\Delta t^2}.}
 \label{eq:rhs}
\end{equation}
The two factors of $\Delta t$ are intentional in this transcription: one
converts a fraction change to a volume rate, and the ordinary projection
scaling supplies the other.

\subsection{Resulting divergence constraint}

Let $u_f$ be velocity in the globally oriented face-normal direction, $A_f$
the face area, and $s_{cf}=\pm1$ convert that direction to the outward normal
of cell $c$. Define integrated divergence
\begin{equation}
 (D u)_c = Q_c = \sum_{f\in\partial c}s_{cf}A_fu_f.
\end{equation}
Positive $Q_c$ denotes net outflow. In the code, the initial flux sum for the
projection right-hand side is $-Q_c^*$.

Write the pressure-to-face-normal-gradient operator as $G$, and let
$R=\operatorname{diag}(1/\rho_f)$. With known pressure boundary contributions
accounted for, the pressure-correction matrix has the form
\begin{equation}
 A=-D R G,\qquad
 A\,\delta p=-\frac{D u^*}{\Delta t}
              -\frac{Vq}{\Delta t^2},
 \label{eq:projection}
\end{equation}
where $Vq$ denotes the vector with entries $V_cq_c$. On an ordinary
interior face the matrix coefficient is
\begin{equation}
 a_f=\frac{\boldsymbol d_f\cdot\boldsymbol N_f}{\rho_f},
\end{equation}
using \code{flow_gradient_coefficients()} for $\boldsymbol d_f$ and the
area-weighted normal $\boldsymbol N_f$. The collapse option changes the
right-hand side, not these coefficients, for a fixed property state.

The ordinary velocity correction remains
\begin{equation}
 u^{\rm new}=u^*-\Delta t\,RG\delta p.
\end{equation}
Combining this with equation~\eqref{eq:projection} gives
\begin{equation}
 \boxed{\frac{Q_c^{\rm new}}{V_c}=-\frac{q_c}{\Delta t}.}
 \label{eq:divergence}
\end{equation}
Thus eligible cells acquire a prescribed net inflow. Elsewhere the usual
zero-divergence target remains. These identities assume consistent discrete
operators and boundary treatment and an exact linear solve; linear-solver
residuals introduce corresponding errors. Ordinary cell-centered velocity
and pressure corrections also follow the projection~\cite{projection}.

\section{How this changes the amount of VOID}

\subsection{Position in the time-step sequence}

The mainline cycle advances material transport before updating flow. With
$n$ denoting the last accepted state, the relevant sequence is~\cite{cycle,flow}
\begin{enumerate}
\item Advect materials using the previously accepted face velocities.
  Apply geometric-tracker boundedness and normalization corrections.
\item If thermal transport is active, update temperature and phase change;
  refresh the phase fractions supplied to flow.
\item Update flow properties and cell classifications. Compare the current
  VOID fraction to the stored accepted fraction to form
  $\delta\alpha_c$ in equation~\eqref{eq:delta}.
\item Advance momentum and project the face velocities with the extra
  source in equation~\eqref{eq:rhs}.
\item Accept the flow/property state. The resulting face velocities drive
  the \emph{next} material-transport update.
\end{enumerate}

The response is therefore delayed through the split sequence. The projection
does not revisit the material update that supplied its current VOID fractions.
Initialization sets the stored fractions to the initial fractions and
explicitly sets \code{void_delta_cc} to zero, so the collapse source is
initially zero~\cite{props}.

\subsection{Preferential removal of VOID by the geometric tracker}

For outward signed material flux volumes $\Phi_{m,cf}$, the tracker first
accumulates fractions according to
\begin{equation}
 \widetilde f_{m,c}=f_{m,c}^{\rm start}
       -\frac{1}{V_c}\sum_f\Phi_{m,cf}.
\end{equation}
It then applies \code{enforce_bounded_vof} after each geometric subcycle.
Following individual fraction cutoff/clipping operations, let
\begin{equation}
 E_c=\sum_m\widetilde f_{m,c}-1,\qquad
 a_c=\sum_{\rm VOID}\widetilde f_{m,c}.
\end{equation}
If $E_c>0$ and $a_c\ge E_c$, it removes the excess solely from VOID:
\begin{equation}
 f_{{\rm VOID},c}^{\rm corrected}
 =\widetilde f_{{\rm VOID},c}\left(1-\frac{E_c}{a_c}\right).
 \label{eq:void-removal}
\end{equation}
With the single VOID slot this is simply $a_c\mapsto a_c-E_c$. In this
branch it leaves real-fluid fractions and their accumulated flux volumes
unchanged. If VOID is insufficient, it removes all available VOID and uses
\code{adjust_flux_all} for the remaining excess. Underfill also reaches the
general adjustment path~\cite{tracker}.

Equation~\eqref{eq:divergence} promotes overfilling, and
equation~\eqref{eq:void-removal} provides the associated VOID loss. For
illustration, if material fluxes sum to the geometric swept volumes, an
initially normalized cell transported for a subsequent interval
$\Delta t_{\rm next}$ with these fixed face velocities has raw overfill
\begin{equation}
 E_c=-\frac{\Delta t_{\rm next}}{V_c}Q_c^{\rm new}
     =q_c\frac{\Delta t_{\rm next}}{\Delta t}.
 \label{eq:next-overfill}
\end{equation}
This is $q_c$ for equal step sizes. Geometry, flux corrections, material
cutoffs, solid occupancy, and exhaustion of VOID can complicate the actual
update. Equation~\eqref{eq:next-overfill} is an interpretation of the
coupling, not a separate formula executed by the tracker.

The preferential VOID correction is present in geometric volume tracking
regardless of the projection's \code{void_collapse} flag. Wisp
redistribution is another, separately controlled operation. The simple
donor-cell tracker does not implement this geometric tracker's boundedness
and preferential-VOID-removal routine; it should not be assumed to provide
the same collapse behavior merely because the projection flag is enabled
~\cite{tracker,simple}.

\section{Consequences of the relaxation rule}

\subsection{Both signs drive contraction; an unchanged fraction does not}

For $\alpha_c=0.02$, $\lambda=0.1$, and $\Delta t=0.01$ seconds, the cap is
$\lambda\alpha_c=0.002$:
\begin{center}
\begin{tabular}{@{}rrr@{}}
\toprule
$\delta\alpha_c$ & $q_c$ & $Q_c^{\rm new}/V_c$ ($\mathrm{s}^{-1}$)\\
\midrule
$0$ & $0$ & $0$\\
$+0.001$ & $0.001$ & $-0.1$\\
$-0.001$ & $0.001$ & $-0.1$\\
$+0.01$ or $-0.01$ & $0.002$ & $-0.2$\\
\bottomrule
\end{tabular}
\end{center}
These are arithmetic examples, not run results. A stationary mixed pocket
with unchanged fraction has no collapse source. Either an increase or a
decrease in its fraction can activate the response. Even translation of a
fragment through cells can produce the Eulerian fraction change used by the
model. Pressure and inertia affect the preceding flow evolution, but the
source contains no explicit pressure or inertial-force threshold.

\subsection{The relaxation parameter is a cap, not a decay prescription}

The model does not directly impose
$\alpha^{n+1}=(1-\lambda)\alpha^n$. The manual's description of roughly
$1/\lambda$ steps to collapse should be read as qualitative guidance.
Only under additional assumptions---persistent saturation of the cap,
equal step sizes, adequate VOID, and no other transport effects---would
repeated removal resemble geometric decay. At $\lambda=0.1$, that idealized
law would leave $0.9^{10}\simeq0.35$ of the original fraction after ten
updates. It would not remove all VOID in ten steps. At $\lambda=1$, complete
removal is permitted by the cap but is not automatically requested, because
$|\delta\alpha_c|$ remains a second limit.

The control is dimensionless and tied to time steps, rather than a fixed
physical relaxation time. At fixed fractions, the added divergence scales
as $1/\Delta t$. Because the resulting velocity is used in the next
transport update, changing step sizes also changes the effective transported
volume as in equation~\eqref{eq:next-overfill}. The relaxation cap is thus
not an unconditional bound on the amount removed in that next update.

\section{Input controls, conservation, and limitations}

The model is selected in the mainline \code{FLOW} namelist:
\begin{lstlisting}
&FLOW
  void_collapse = T
  void_collapse_relaxation = 0.1
/
\end{lstlisting}
The default is disabled. The relaxation factor defaults to $0.1$, and the
namelist reader validates $0\le\lambda\le1$. The bridging parameter list
uses \code{options/void-collapse} and
\code{options/void-collapse-relaxation}. The projection applies the default
relaxation value when it reads that list. Setting $\lambda=0$ gives a zero
additional source~\cite{input,projection,manual}.

\textbf{Conservation.} VOID volume is deliberately not conserved by its
normalization correction. Removing VOID alone from an overfilled cell need
not change real-fluid mass. However, geometric tracking also uses cutoffs
and cell-local flux adjustments. The source explicitly notes that those
adjustments can make material flux volumes inconsistent across a shared
face. Hence exact global real-fluid conservation is not guaranteed by the
complete algorithm. The manual reports increased mass loss with larger
relaxation factors (of order $0.5\%$ in its experience) and improvement with
more subcycling; those statements are empirical guidance, not error bounds
or measurements made for this note~\cite{tracker,manual}.

\textbf{Global compatibility.} The collapse source is local and has no
global compensation mechanism. The discrete divergence theorem requires
the sum of integrated divergences in a connected active region to equal
its net boundary flux. Consequently, an entirely closed region with
impermeable boundaries cannot satisfy a strictly negative total prescribed
divergence. A free surface or open boundary can accommodate compensating
motion. With only Neumann pressure conditions, pinning a pressure equation
does not cure the compatibility condition of the original divergence
equations; it permits a compensating defect in the modified equation. This
is a mathematical limitation of applying a net source in that setting, not
a failure reproduced in a run here.

\textbf{Geometry and resolution.} Eligibility depends on face adjacency
and the existence of pure void cells. Refinement or mesh perturbation can
therefore alter eligibility, even for similar physical interfaces. The
model does not measure cavity size, reconstruct a closed void surface, or
determine connectivity across the domain. Its local criterion should not
be equated with identifying every physically trapped bubble.

\textbf{Current test evidence.} The repository's \code{test/void} example
constructs a restart with an injected VOID fraction and enables collapse.
Its Python check bounds the maximum absolute change of the VOID field from
the initial output by $10^{-4}$. Although its comment speaks of preventing
growth, the comparison bounds changes in either direction. It does not
validate a collapse rate, the prescribed divergence, or liquid mass
conservation~\cite{test}. Source inspection establishes the implemented
rule; it does not establish accuracy or robustness for every geometry.

\section{Source guide and document reproduction}

The links below refer to the immutable mainline revision stated on the
first page. Within a clone, for example, the core routine is available with
\begin{lstlisting}
git show 344a7c3a:src/truchas/flow/flow_projection_type.F90
\end{lstlisting}
To rebuild this document with pdfLaTeX, run twice:
\begin{lstlisting}
pdflatex -interaction=nonstopmode -halt-on-error void-collapse-model.tex
\end{lstlisting}

\begin{thebibliography}{10}
\small
\bibitem{projection}
\srclink{src/truchas/flow/flow_projection_type.F90}{201}{\code{flow_projection_type.F90}}.
\code{init}: defaults; \code{setup_solver}: collapse source, pressure matrix,
and pinning; \code{velocity_fc_correct}: face-velocity correction.

\bibitem{props}
\srclink{src/truchas/flow/flow_props_type.F90}{176}{\code{flow_props_type.F90}}.
\code{set_initial_state}, \code{update_cc}, \code{update_fc}, and
\code{accept}: fraction history, cell/face classification, and state storage.

\bibitem{cycle}
\srclink{src/truchas/drivers/drivers.F90}{218}{\code{drivers.F90}}, main cycle;
\srclink{src/truchas/physics/volume_tracking/vtrack_driver.F90}{235}{\code{vtrack_driver.F90}},
\code{vtrack_update};
\srclink{src/truchas/flow/flow_driver.F90}{336}{\code{flow_driver.F90}},
\code{flow_step} and \code{flow_accept}.

\bibitem{flow}
\srclink{src/truchas/flow/flow_type.F90}{336}{\code{flow_type.F90}}.
\code{step}: property update, prediction, projection; \code{accept}:
velocity history.

\bibitem{tracker}
\srclink{src/truchas/physics/volume_tracking/geometric_volume_tracker_type.F90}{910}{\code{geometric_volume_tracker_type.F90}}.
\code{flux_volumes}, \code{accumulate_volume}, \code{enforce_bounded_vof},
\code{adjust_flux_matl}, and \code{adjust_flux_all}.

\bibitem{simple}
\srclink{src/truchas/physics/volume_tracking/simple_volume_tracker_type.F90}{49}{\code{simple_volume_tracker_type.F90}}.
\code{flux_volumes}: donor-cell update.

\bibitem{input}
\srclink{src/truchas/flow/flow_namelist.F90}{203}{\code{flow_namelist.F90}}.
Namelist defaults, relaxation bounds, and parameter-list translation.

\bibitem{manual}
\srclink{doc/reference-manual/FLOW_Namelist/index.rst}{220}{Truchas reference manual, \code{FLOW} namelist}.
Entries \code{void_collapse} and \code{void_collapse_relaxation}.

\bibitem{test}
\srclink{test/void/README.md}{1}{\code{test/void/README.md}},
\srclink{test/void/add-solid-void.py}{1}{\code{add-solid-void.py}},
\srclink{test/void/test.inp}{52}{\code{test.inp}}, and
\srclink{test/void/test.py}{1}{\code{test.py}}.
Restart preparation, input configuration, and existing test criterion.
\end{thebibliography}
\end{document}
